\section{The Ticket Urgency Score}
\label{sec:score}

\subsection{Definition}

For a ticket $t$ evaluated at decision time $\tau$, the score is a static
component under multiplicative aging:
\begin{align}
U(t,\tau) &= S(t)\,\bigl(1 + \alpha\, a(t,\tau)\bigr) \label{eq:tus}\\
S(t) &= w_P\,\hat{\mathbb{E}}[\mathrm{sev}\mid t] + w_S\,\hat{P}(\mathrm{Neg}\mid t)
        + w_I\,\kappa(\hat{\imath}(t)) \nonumber\\
a(t,\tau) &= \bigl(\tau - \mathrm{arrival}(t)\bigr) / D \nonumber
\end{align}
with $w \ge 0$, $\sum w = 1$, $\alpha \ge 0$, and $D$ the service window of the
ticket's class. Four numbers are fitted. Waiting tickets are served in
decreasing $U$.

\textbf{Expected severity, not the argmax.}
$\hat{\mathbb{E}}[\mathrm{sev}\mid t] = \sum_k \hat{p}_k(t)\,\mathrm{sev}_k$ over
$\mathrm{sev} = (0, \tfrac12, 1)$. Taking the argmax of the priority head and
mapping it to a tier is what a naive deployment does, and it discards exactly
what the queue needs: confidence. A confidently-Low ticket and a $0.49/0.51$
coin flip both become "Low" and then sit in the same tier in arrival order. Our
error analysis finds priority errors concentrate on the Low/Medium boundary,
which is precisely where the argmax destroys the most information. Ordering by
the posterior rather than its argmax is the property Argon and Ziya
\cite{argon2009priority} prove lowers long-run cost, so this is a theorem rather
than a preference.

\textbf{Sentiment as a nudge, never a gate.} At 0.7138 \negf{} the sentiment head
misses roughly a third of the negative tickets. A weak signal carrying a
small fitted weight inside a score is useful; the same signal as a routing rule
would be unsafe.

\textbf{Learned per-intent criticality.} $\kappa(i) = \hat{P}(\text{High} \mid
\text{intent}=i)$, 77 values estimated on training and development data only.
This is the term that makes~\eqref{eq:tus} a $c\mu$ rule \cite{cox1961queues}
with a \emph{learned} cost rather than an assigned one. It ranks
\texttt{compromised\_card}, \texttt{lost\_or\_stolen\_card}, and
\texttt{card\_payment\_not\_recognised} highest, which is the ordering a bank
would choose by hand. It also matters for Section~\ref{sec:disparity}: intent is
the most accurate of the three heads, so $\kappa$ acts as a script-robust prior
that backstops the priority head where the text is degraded.

\textbf{Aging is multiplicative}, so a stale Low climbs but still climbs more
slowly than a stale High, which additive aging cannot express.

\subsection{Fitting}

Weights are fitted by grid search on development data over the 3-simplex at
$0.05$ resolution with a grid on $\alpha$, then applied unchanged to test:
$w = (0.80, 0.10, 0.10)$, $\alpha = 16$. Development posteriors come from a
train-only fit, because the saved encoder checkpoints were fitted on training
plus development and have memorised it.

\textbf{The exact vector is not the finding, and we do not present it as one.}
The surface is flat: 71 of 231 simplex points sit within 5\% of the optimum,
spanning $w_P \in [0.10, 0.90]$. What the fit establishes is that priority
dominates and that no corner of the simplex is competitive, not that $0.80$ beats
$0.70$. We publish the whole surface rather than the argmin. Extending the
$\alpha$ grid to 128 confirms the fitted value sits on a plateau, so aging is a
policy dial rather than a fitted constant.

An ablation at $\rho = 1.05$ shows the terms are not interchangeable. Priority
alone reaches 2.15 minutes mean wait for urgent tickets against the full score's
1.97, and priority with intent reaches 1.85; sentiment on its own degrades the
head of the queue. Intent is what earns its place, and it does so through
$\kappa$ rather than as a third severity estimate.
